\documentclass[11pt]{article}

\usepackage[margin=1in]{geometry}
\usepackage{amsmath,amssymb,amsthm,mathtools}
\usepackage{booktabs}
\usepackage{enumitem}
\usepackage{microtype}
\usepackage[hidelinks]{hyperref}

\newcommand{\R}{\mathbb R}
\newcommand{\C}{\mathbb C}
\newcommand{\N}{\mathbb N}
\newcommand{\Nzero}{\mathbb N_0}
\newcommand{\D}{\mathbb D}
\newcommand{\E}{\mathbb E}
\newcommand{\Prob}{\mathbb P}
\newcommand{\1}{\mathbf 1}
\newcommand{\sgn}{\operatorname{sgn}}
\newcommand{\Hol}{\operatorname{Hol}}
\newcommand{\dist}{\operatorname{dist}}
\newcommand{\ol}{\overline}
\newcommand{\ip}[2]{\langle #1,#2\rangle}
\newcommand{\KG}{K_G}
\newcommand{\SDP}{\operatorname{SDP}}
\newcommand{\OPT}{\operatorname{OPT}}
\newcommand{\Khyp}{K_{\mathrm{hyp}}}
\newcommand{\rhohyp}{\rho_{\mathrm{hyp}}}
\newcommand{\gammazero}{\gamma_0}
\newcommand{\eps}{\varepsilon}
\newcommand{\calA}{\mathcal A}
\newcommand{\calH}{\mathcal H}
\newcommand{\calK}{\mathcal K}
\DeclareMathOperator{\arsinh}{arsinh}

\theoremstyle{plain}
\newtheorem{theorem}{Theorem}[section]
\newtheorem{lemma}[theorem]{Lemma}
\newtheorem{proposition}[theorem]{Proposition}
\newtheorem{corollary}[theorem]{Corollary}

\theoremstyle{definition}
\newtheorem{definition}[theorem]{Definition}
\newtheorem{assumption}[theorem]{Certificate Module}

\theoremstyle{remark}
\newtheorem{remark}[theorem]{Remark}

\title{A Gaussianized Chaos Krivine Scheme for the Grothendieck Constant}
\author{Draft}
\date{June 2026}

\begin{document}
\maketitle

\begin{abstract}
We introduce and analyze a family of high-dimensional Krivine rounding schemes built from Gaussianized Hermite-chaos reservoirs.  The simplest member of the family uses an anti-aligned third-chaos reservoir and an aligned fifth-chaos reservoir.  Its finite-dimensional incarnations are genuine Gaussian rounding schemes, while its limiting correlation model remembers the reservoir degrees through the weighted correlations \(t^3\) and \(t^5\).  We derive an explicit coefficient formula for the limiting correlation function in terms of two-dimensional Gaussian integrals, and we prove a finite-dimensional transfer theorem: any strict inverse-majorant certificate for the limiting model transfers to all sufficiently high finite-dimensional schemes.  The transfer theorem is analytic rather than coefficientwise; it uses locally uniform convergence, stability of inverse branches by Rouche's theorem, and Cauchy estimates for the inverse-tail.

For the parameters
\[
  \eta=0.136419125,\qquad s_3=0.34101124,\qquad s_5=0.05276111,
\]
the limiting normalized correlation has linear coefficient \(b_1=0.8815738220\ldots\), while the cubic and quintic coefficients are numerically annihilated.  We reduce the proof of admissibility at
\[
  \gammazero=0.8815624
\]
to a short interval-arithmetic certificate for the head and tail of the limiting coefficient sequence.  Once this certificate is inserted, the construction gives
\[
  \KG \le {\pi\over 2\gammazero}
  =1.7818322637\ldots
  ={\pi\over 2\log(1+\sqrt2)}-3.81714\cdot 10^{-4}\ldots .
\]
The certificate section of this draft is deliberately modular: all preceding arguments are analytic and exact, and the remaining computer-assisted inequalities are isolated explicitly.
\end{abstract}


\section{Introduction}

The Grothendieck inequality~\cite{Grothendieck1953} from functional analysis shows that~\cite{LindenstraussPelczynski1968}, there exists a universal constant $K<\infty$ such that for every $m,n\in\N$ and every matrix $A=(a_{ij})\in\R^{m\times n}$,
\begin{equation}\label{eq:grothendieck-ineq}
  \max_{\eps_i,\delta_j\in\{\pm1\}}
  \left|\sum_{i=1}^m\sum_{j=1}^n a_{ij}\eps_i\delta_j\right|
  \le
  \sup_{u_i,v_j\in S^{m+n-1}}
  \left|\sum_{i=1}^m\sum_{j=1}^n a_{ij}\ip{u_i}{v_j}\right|
  \le
  K\max_{\eps_i,\delta_j\in\{\pm1\}}
  \left|\sum_{i=1}^m\sum_{j=1}^n a_{ij}\eps_i\delta_j\right|.
\end{equation}
Here $S^{m+n-1}$ denotes the unit sphere in $\R^{m+n}$. The Grothendieck constant $\KG$ is the infimum of $K$ for which this inequality holds.

The Grothendieck inequality is equivalent to the statement that the integrality gap of the canonical SDP relaxation for maximizing the bilinear form $x^\top Ay$ over $x\in\{-1,1\}^m$ and $y\in\{-1,1\}^n$ is an absolute constant $\KG$. As a result, one obtains efficient algorithms for approximating the matrix cut norms that, among other applications, lead to algorithms for finding Szemeredi partitions of graphs~\cite{AlonNaor2006}. Beyond this natural role in the design of approximation algorithms in computer science, the inequality is important in several subfields, including the geometry of Banach spaces, $C^*$-algebras, harmonic analysis, operator spaces, and quantum mechanics; see, for example,~\cite{LindenstraussTzafriri2013,Pisier2012,KhotNaor2011}.

Despite a long line of work on the problem, the exact value of $\KG$ remains unknown. Grothendieck proved that $\KG\le \sinh(\pi/2)=2.3012\ldots$. The next advance was made in 1977 by Krivine~\cite{Krivine1977}, who proved that
\begin{equation}\label{eq:krivine-bound-intro}
  \KG\le {\pi\over 2\log(1+\sqrt2)}=1.7822\ldots .
\end{equation}
Krivine's proof involved developing an explicit \emph{rounding scheme}, an algorithm that maps an input set of vectors $u_i,v_j\in S^{d-1}$ into signs $\eps_i,\delta_j\in\{\pm1\}$. Krivine's scheme consisted of first lifting $\{u_i,v_j\}$ to an infinite-dimensional space and then partitioning them by a random hyperplane into signed $\{\eps_i,\delta_j\}$. He showed that this hyperplane partition yields the above bound and conjectured that it is optimal, leading to subsequent efforts to find matching lower bounds~\cite{Davie1984,Reeds1991}.

Krivine's conjecture was refuted three decades later by Braverman, Makarychev, Makarychev, and Naor~\cite{BMMN2011}, who showed that, in fact,
\[
  \KG < {\pi\over 2\log(1+\sqrt2)}-\eps
\]
for some $\eps>0$. Their proof relied on a ``mixing'' of two rounding strategies, i.e., applying a new non-hyperplane scheme with a small probability and the hyperplane rounding scheme with the rest. This idea of mixing is natural in the context of analyzing rounding algorithms and is informally analogous to first-order methods that obtain a better guess for an optimum by exploiting a nonzero gradient. Such a first-order analysis suffices to prove the existence of a scheme that outperforms Krivine's rounding. However, this technique makes it difficult to derive an explicit improvement.

In the present work, we develop a different high-dimensional rounding mechanism.  Instead of mixing a small low-dimensional perturbation with the hyperplane scheme, we build finite-dimensional schemes from Hermite-chaos reservoirs and analyze their limiting weighted Gaussian correlation model.  The resulting certificate, once the interval-arithmetic module stated below is filled in, gives an explicit improvement over the hyperplane value.

\subsection{The result in certificate-reduction form}

Write
\[
  \rhohyp:=\log(1+\sqrt2),
  \qquad
  \Khyp:={\pi\over2\rhohyp}={\pi\over2\log(1+\sqrt2)}.
\]
The numerical target in this paper is
\[
  \gammazero=0.8815624>\rhohyp.
\]
This corresponds to the raw correlation value
\[
  c_0={2\gammazero\over \pi}=0.5612200543\ldots
\]
and to the Grothendieck bound
\[
  {1\over c_0}={\pi\over 2\gammazero}=1.7818322637\ldots .
\]
The improvement over the hyperplane/Krivine value is
\[
  \Khyp-{\pi\over 2\gammazero}=3.817144799\cdot 10^{-4}\ldots .
\]

\begin{theorem}[Main theorem, conditional on the coefficient certificate]\label{thm:main}
Assume the coefficient certificate stated in Certificate Module~\ref{ass:certificate}.  Then there is a finite-dimensional Krivine rounding scheme in the Gaussianized third/fifth-chaos family such that
\[
  \KG \le {\pi\over 2\gammazero}
  < {\pi\over 2\log(1+\sqrt2)}-3.81714\cdot 10^{-4}.
\]
Equivalently,
\[
  \KG \le 1.7818322637\ldots .
\]
\end{theorem}

The conditional phrasing is intentional.  The analytic part of the proof is complete in this draft.  What remains is a finite interval-arithmetic certificate for the explicit coefficient inequalities in Certificate Module~\ref{ass:certificate}.  Once that module is replaced by an Arb or comparable outward-rounded certificate, Theorem~\ref{thm:main} becomes unconditional.

\subsection{Overview of the construction}

Let \(H_n\) denote the orthonormal probabilists' Hermite polynomial.  For every reservoir length \(N\), define
\[
  P_{3,N}=N^{-1/2}\sum_{j=1}^N H_3(U_j),
  \qquad
  P_{5,N}=N^{-1/2}\sum_{j=1}^N H_5(V_j),
\]
where the variables are independent standard Gaussians.  The finite rounding pair is
\begin{align*}
  f_N &= \sgn\bigl(Z+\eta H_3(X)+s_3P_{3,N}+s_5P_{5,N}\bigr),\\
  g_N &= \sgn\bigl(Z-\eta H_3(X)-s_3P_{3,N}+s_5P_{5,N}\bigr).
\end{align*}
Thus degree \(3\) is anti-aligned between the two sides, while degree \(5\) is aligned.  The sign pattern is not cosmetic: perturbatively, odd degree \(2k+1\) naturally wants relative sign \((-1)^k\), so degrees \(3\) and \(5\) want opposite and equal signs, respectively.

For fixed \(t\in(-1,1)\), the paired reservoirs have asymptotic covariance
\[
  \E[P_{3,N}P_{3,N}']=t^3,
  \qquad
  \E[P_{5,N}P_{5,N}']=t^5.
\]
Hence the limiting model is a four-dimensional Gaussian threshold pair
\begin{align*}
  f(z,x,r_3,r_5)&=\sgn\bigl(z+\eta H_3(x)+s_3r_3+s_5r_5\bigr),\\
  g(z,x,r_3,r_5)&=\sgn\bigl(z-\eta H_3(x)-s_3r_3+s_5r_5\bigr),
\end{align*}
but with coordinate correlations
\[
  \rho_Z=t,\quad \rho_X=t,
  \qquad \rho_3=t^3,
  \qquad \rho_5=t^5.
\]
This weighted correlation is the right limiting object for the finite pure-chaos schemes.

The limiting normalized correlation is an odd series
\[
  H(t)=\sum_{m\ge1}b_m t^m.
\]
For the displayed parameters, the leading coefficients are numerically
\[
  b_1=0.8815738220496\ldots,
  \qquad b_3=3.45\cdot 10^{-9},
  \qquad b_5=-9.04\cdot 10^{-11}.
\]
Since the inverse coefficients satisfy
\[
  A_1={1\over b_1},\qquad
  A_3=-{b_3\over b_1^4},\qquad
  A_5={3b_3^2-b_1b_5\over b_1^7},
\]
the annihilation of \(b_3\) and \(b_5\) removes the first two nonlinear inverse obstructions.



\subsection{Proof architecture}\label{subsec:architecture}

The proof is organized so that every analytic reduction is separated from the eventual interval-arithmetic certificate.  The section headers below match the rest of the paper.

\paragraph{Section~\ref{sec:krivine}: Krivine preprocessing.}
This section states the inverse-majorant interface.  A normalized correlation function $H$ is admissible at radius $\gamma$ if its local inverse $G=H^{-1}=\sum A_n\zeta^n$ is holomorphic near $\gamma\overline{\D}$ and satisfies
\[
  \sum_{n\ge1}|A_n|\gamma^n\le 1.
\]
The conclusion is $\KG\le \pi/(2\gamma)$.  The proof is standard and is placed in Appendix~\ref{app:preprocessing}.

\paragraph{Section~\ref{sec:scheme}: The finite family and the limiting model.}
Definitions~\ref{def:finite-family} and~\ref{def:finite-35} define genuine finite-dimensional Gaussian rounding schemes built from normalized Hermite reservoirs.  Definition~\ref{def:limiting-model} defines the limiting weighted Gaussian model, where the third and fifth reservoirs retain the correlation powers $t^3$ and $t^5$.

\paragraph{Section~\ref{sec:coefficients}: Coefficients of the limiting model.}
Definitions~\ref{def:reservoir-normalization}, \ref{def:qk}, and~\ref{def:Cabk} reduce the limiting threshold calculation to explicit one- and two-dimensional Gaussian quantities.  Proposition~\ref{prop:hermite-coefficients} gives the Hermite coefficients of the limiting thresholds, and Theorem~\ref{thm:coefficient-formula} converts them into the weighted-degree formula for $b_m=[t^m]H(t)$.

\paragraph{Section~\ref{sec:transfer}: Finite-dimensional transfer.}
This is the rigorously delicate part.  Proposition~\ref{prop:loc-unif} proves locally uniform convergence $H_N\to H$ on the unit disk.  Theorem~\ref{thm:finite-transfer} then transfers a strict inverse-majorant certificate from $H$ to $H_N$ using Rouche's theorem and Cauchy estimates.  Fixed-degree coefficient convergence alone would not be enough, because it does not control the inverse tail.

\paragraph{Section~\ref{sec:wiener}: The limiting inverse certificate.}
Definition~\ref{def:wiener-data} introduces the Wiener-algebra quantities used to control the inverse.  Theorem~\ref{thm:aposteriori} gives an a posteriori inverse-tail certificate from a short inverse head and a small coefficient tail for $H(t)=b_1t+E(t)$.

\paragraph{Section~\ref{sec:certification}: Head/tail certification and proof of Theorem~\ref{thm:main}.}
Certificate Module~\ref{ass:certificate} states the finite inequalities that remain to be checked by interval arithmetic.  Proposition~\ref{prop:lim-admissible} plugs those inequalities into Theorem~\ref{thm:aposteriori}, and the proof of Theorem~\ref{thm:main} combines this with the finite-dimensional transfer theorem and Krivine preprocessing.

\paragraph{Section~\ref{sec:discussion}: Discussion and future directions.}
The final section explains why the third/fifth sign pattern cancels the first inverse obstructions, records a possible $7/9$ extension, and discusses how one could extract effective finite dimensions from the collision expansion.

\section{Krivine preprocessing}\label{sec:krivine}

This short section records the analytic interface between a correlation function and an upper bound on the Grothendieck constant.  It is the first step in the proof architecture from Section~\ref{subsec:architecture}; the full preprocessing proof, including the mixed version used in earlier works, is deferred to Appendix~\ref{app:preprocessing}.

Let $f,g:\R^k\to\{-1,1\}$ be odd measurable functions and let
\[
  H_{f,g}(t)={\pi\over2}\E\left[f(G_1)g\left(tG_1+\sqrt{1-t^2}G_2\right)\right]
\]
be the arcsine-normalized Gaussian correlation, where $G_1,G_2$ are independent standard Gaussian vectors in $\R^k$.  The preprocessing criterion says that if the local inverse
\[
  H_{f,g}^{-1}(\zeta)=\sum_{n\ge1}A_n\zeta^n
\]
is holomorphic near $\gamma\overline{\D}$ and satisfies
\[
  \sum_{n\ge1}|A_n|\gamma^n\le1,
\]
then
\[
  \KG\le {\pi\over2\gamma}.
\]
This is Theorem~\ref{thm:krivine-criterion} in Appendix~\ref{app:preprocessing}.  For the hyperplane scheme $f=g=\sgn$, one has $H_{f,g}(t)=\arcsin t$ and $H_{f,g}^{-1}(\zeta)=\sin\zeta$, so the largest radius detected by this criterion is $\rho_{\mathrm{hyp}}=\log(1+\sqrt2)$.

The rest of the paper is devoted to constructing a finite-dimensional scheme whose limiting correlation function is admissible at a radius larger than $\rho_{\mathrm{hyp}}$, and then proving that the limiting certificate transfers back to sufficiently high finite dimensions.

\section{The Gaussianized chaos family}\label{sec:scheme}

This section is Section~\ref{sec:scheme} in the proof architecture of Section~\ref{subsec:architecture}.  We first define the genuine finite-dimensional rounding schemes, and then define the limiting weighted model that captures their fixed-degree asymptotics.

\subsection{Hermite normalization}

Throughout the paper, \(H_n\) denotes the orthonormal probabilists' Hermite polynomial, so that
\[
  \E[H_m(Z)H_n(Z)]=\delta_{mn},\qquad Z\sim N(0,1).
\]
Thus
\[
  H_0=1,
  \qquad H_1(x)=x,
  \qquad H_3(x)={x^3-3x\over\sqrt6},
  \qquad H_5(x)={x^5-10x^3+15x\over\sqrt{120}}.
\]
If \((X,Y)\) are standard Gaussians with correlation \(t\), then
\begin{equation}\label{eq:hermite-cov}
  \E[H_m(X)H_n(Y)]=\delta_{mn}t^n.
\end{equation}


\subsection{Finite-dimensional schemes}\label{subsec:finite-schemes}

\begin{definition}[Finite Gaussianized-chaos family]\label{def:finite-family}
Let $D\subseteq\{3,5,7,9,\ldots\}$ be a finite set of odd reservoir degrees.  For each $d\in D$, fix a real weight $s_d$, and define the relative sign
\[
  \sigma_d=(-1)^{(d-1)/2}.
\]
Thus $\sigma_3=-1$, $\sigma_5=+1$, $\sigma_7=-1$, and so on.

For reservoir length $N$, let
\[
  P_{d,N}=N^{-1/2}\sum_{j=1}^N H_d(U_{d,j}),
\]
where all variables $Z,X,U_{d,j}$ are independent standard Gaussians.  The finite rounding functions are
\begin{align}
  f_N&=\sgn\biggl(Z+\eta H_3(X)+\sum_{d\in D}s_dP_{d,N}\biggr),\label{eq:finite-f-general}\\
  g_N&=\sgn\biggl(Z-\eta H_3(X)+\sum_{d\in D}\sigma_d s_dP_{d,N}\biggr).\label{eq:finite-g-general}
\end{align}
They are odd measurable functions on $\R^{2+|D|N}$.  The value of $\sgn(0)$ is immaterial, since every threshold event has probability zero.
\end{definition}

\begin{definition}[The third/fifth finite scheme]\label{def:finite-35}
The scheme used for the numerical certificate is the specialization of Definition~\ref{def:finite-family} with $D=\{3,5\}$ and
\begin{equation}\label{eq:params}
  \eta=0.136419125,
  \qquad
  s_3=0.34101124,
  \qquad
  s_5=0.05276111.
\end{equation}
Equivalently,
\begin{align}
  f_N&=\sgn\bigl(Z+\eta H_3(X)+s_3P_{3,N}+s_5P_{5,N}\bigr),\label{eq:finite-f-35}\\
  g_N&=\sgn\bigl(Z-\eta H_3(X)-s_3P_{3,N}+s_5P_{5,N}\bigr).\label{eq:finite-g-35}
\end{align}
Thus degree $3$ is anti-aligned between the two sides, while degree $5$ is aligned.
\end{definition}


\subsection{The limiting weighted model}\label{subsec:limiting-model}

\begin{definition}[Limiting weighted Gaussian model]\label{def:limiting-model}
Let $D,\eta,(s_d)_{d\in D}$ be as in Definition~\ref{def:finite-family}.  Let $R_d$, $d\in D$, be independent standard Gaussians.  The limiting threshold pair is
\begin{align}
  f(z,x,(r_d))&=\sgn\biggl(z+\eta H_3(x)+\sum_{d\in D}s_dr_d\biggr),\label{eq:lim-f-general}\\
  g(z,x,(r_d))&=\sgn\biggl(z-\eta H_3(x)+\sum_{d\in D}\sigma_d s_dr_d\biggr).\label{eq:lim-g-general}
\end{align}
For the limiting correlation at parameter $t$, the paired coordinates satisfy
\[
  \rho_Z=t,
  \qquad \rho_X=t,
  \qquad \rho_d=t^d\quad(d\in D),
\]
and different coordinate pairs are independent.  The arcsine-normalized limiting correlation is
\begin{equation}\label{eq:lim-H}
  H(t)={\pi\over2}\E\bigl[f(Z,X,(R_d))g(Z',X',(R_d'))\bigr].
\end{equation}
\end{definition}

Because all degrees in the model are odd, $H$ is an odd function.  For $D=\{3,5\}$, Definition~\ref{def:limiting-model} says that the reservoir coordinates have correlations $t^3$ and $t^5$.  This is the formal memory of the fact that the finite reservoirs came from pure third and fifth Gaussian chaoses.


\section{The coefficient formula for the limiting model}\label{sec:coefficients}

This section implements the coefficient step in the proof architecture, namely Section~\ref{sec:coefficients} of the roadmap in Section~\ref{subsec:architecture}.  The point is to reduce the four-dimensional limiting threshold model, with weighted correlations, to explicit two-dimensional Gaussian integrals.

\subsection{Reservoir-direction data}\label{subsec:reservoir-data}

\begin{definition}[Reservoir normalization]\label{def:reservoir-normalization}
For a finite reservoir set $D$ and weights $(s_d)_{d\in D}$, define
\[
  S=\biggl(\sum_{d\in D}s_d^2\biggr)^{1/2},
  \qquad
  \alpha_d={s_d\over S}\quad(d\in D).
\]
Thus $\sum_{d\in D}\alpha_d^2=1$, and
\[
  \sum_{d\in D}s_dr_d=S\sum_{d\in D}\alpha_dr_d.
\]
\end{definition}

\begin{definition}[Reservoir threshold coefficients]\label{def:qk}
Let $R\sim N(0,1)$.  For $u\in\R$ and $k\in\Nzero$, define
\[
  q_k(u)=\E_R[\sgn(u+R)H_k(R)].
\]
These are the Hermite coefficients of the one-dimensional threshold $R\mapsto \sgn(u+R)$.
\end{definition}

\begin{lemma}[Closed form for $q_k$]\label{lem:qk}
For all $u\in\R$,
\begin{align}
  q_0(u)&=\operatorname{erf}\left({u\over\sqrt2}\right),\label{eq:q0}\\
  q_k(u)&=2\varphi(u){H_{k-1}(-u)\over\sqrt{k}},\qquad k\ge1,\label{eq:qk}
\end{align}
where $\varphi(u)=(2\pi)^{-1/2}e^{-u^2/2}$.
\end{lemma}

\begin{proof}
For $k=0$,
\[
  q_0(u)=\Prob[R>-u]-\Prob[R<-u]=\operatorname{erf}\left({u\over\sqrt2}\right).
\]
For $k\ge1$, orthogonality to constants gives
\[
  q_k(u)=2\int_{-u}^{\infty}H_k(r)\varphi(r)\,dr.
\]
The orthonormal Hermite polynomials satisfy
\[
  {d\over dr}\left(\varphi(r)H_{k-1}(r)\right)
  =-\sqrt{k}\,\varphi(r)H_k(r),
\]
so integration from $-u$ to $\infty$ gives \eqref{eq:qk}.
\end{proof}

\begin{definition}[Two-dimensional coefficient integrals]\label{def:Cabk}
For $a,b,k\in\Nzero$, define
\begin{equation}\label{eq:C-abk}
  C_{a,b,k}
  =\E_{Z,X}\left[H_a(Z)H_b(X)
    q_k\left({Z+\eta H_3(X)\over S}\right)\right],
\end{equation}
where $Z,X$ are independent standard Gaussians.  Equivalently, $C_{a,b,k}$ is the two-dimensional integral
\[
  \int_{\R^2}H_a(z)H_b(x)
  q_k\left({z+\eta H_3(x)\over S}\right)\varphi(z)\varphi(x)\,dz\,dx.
\]
\end{definition}

The quantities $C_{a,b,k}$ are the only continuous integrals appearing in the coefficient formula.  Everything else is finite weighted-degree bookkeeping.

\subsection{Hermite coefficients of the limiting thresholds}\label{subsec:threshold-hermite-coefficients}

For a multi-index $\beta=(\beta_d)_{d\in D}\in\Nzero^D$, write
\[
  |\beta|=\sum_{d\in D}\beta_d,
  \qquad
  \beta!=\prod_{d\in D}\beta_d!,
  \qquad
  \alpha^\beta=\prod_{d\in D}\alpha_d^{\beta_d}.
\]

\begin{proposition}[Hermite coefficients of the limiting thresholds]\label{prop:hermite-coefficients}
Let $f,g$ be the limiting threshold functions in Definition~\ref{def:limiting-model}.  Write their Hermite expansions in the tensor basis as
\[
  f=\sum_{a,b,\beta} A_{a,b,\beta}\,H_a(Z)H_b(X)\prod_{d\in D}H_{\beta_d}(R_d),
\]
\[
  g=\sum_{a,b,\beta} B_{a,b,\beta}\,H_a(Z)H_b(X)\prod_{d\in D}H_{\beta_d}(R_d).
\]
Then, for $k=|\beta|$,
\begin{equation}\label{eq:f-coeff}
  A_{a,b,\beta}
  =C_{a,b,k}\sqrt{k!\over \beta!}\,\alpha^\beta,
\end{equation}
while
\begin{equation}\label{eq:g-coeff}
  B_{a,b,\beta}
  =(-1)^b\prod_{d\in D}\sigma_d^{\beta_d}\,A_{a,b,\beta}.
\end{equation}
\end{proposition}

\begin{proof}
Set
\[
  T(Z,X)=Z+\eta H_3(X),
  \qquad
  R_\alpha=\sum_{d\in D}\alpha_dR_d.
\]
Then
\[
  f=\sgn(T+SR_\alpha).
\]
Conditioning on $Z,X$, the coefficient of $H_k(R_\alpha)$ in the $R_\alpha$ direction is
\[
  q_k\left({T(Z,X)\over S}\right).
\]
Taking the additional inner product against $H_a(Z)H_b(X)$ gives $C_{a,b,k}$.

The multivariate Hermite addition formula gives
\begin{equation}\label{eq:hermite-addition}
  H_k\left(\sum_{d\in D}\alpha_dR_d\right)
  =\sum_{\substack{\beta\in\Nzero^D\\ |\beta|=k}}
  \sqrt{k!\over \beta!}\,\alpha^\beta
  \prod_{d\in D}H_{\beta_d}(R_d).
\end{equation}
Combining the previous paragraph with \eqref{eq:hermite-addition} proves \eqref{eq:f-coeff}.

Finally,
\[
  g(z,x,(r_d))=f(z,-x,(\sigma_dr_d)_{d\in D}).
\]
Since $H_b(-x)=(-1)^bH_b(x)$ and $H_{\beta_d}(\sigma_dr_d)=\sigma_d^{\beta_d}H_{\beta_d}(r_d)$, the coefficient relation \eqref{eq:g-coeff} follows.
\end{proof}

\subsection{The weighted-degree coefficient formula}\label{subsec:weighted-degree-formula}

\begin{theorem}[Weighted-degree coefficient formula]\label{thm:coefficient-formula}
Let
\[
  H(t)=\sum_{m\ge1}b_m t^m
\]
be the limiting normalized correlation \eqref{eq:lim-H}.  Then
\begin{equation}\label{eq:general-bm}
  b_m={\pi\over2}
  \sum_{\substack{a,b\in\Nzero,\ \beta\in\Nzero^D\\ a+b+\sum_{d\in D}d\beta_d=m}}
  (-1)^b\prod_{d\in D}\sigma_d^{\beta_d}
  {k!\over \beta!}
  \prod_{d\in D}\alpha_d^{2\beta_d}
  C_{a,b,k}^2,
  \qquad k=|\beta|.
\end{equation}
Only odd $m$ occur.
\end{theorem}

\begin{proof}
Under the weighted correlation model in Definition~\ref{def:limiting-model}, the Hermite covariance identity \eqref{eq:hermite-cov} contributes
\[
  t^a t^b\prod_{d\in D}(t^d)^{\beta_d}
  =t^{a+b+\sum_d d\beta_d}
\]
to the tensor basis element indexed by $(a,b,\beta)$.  Multiplying the $f$- and $g$-coefficients from Proposition~\ref{prop:hermite-coefficients} and summing over all indices of weighted degree $m$ gives \eqref{eq:general-bm}.

The functions $f$ and $g$ are odd under simultaneous negation of all coordinates.  Hence their Hermite coefficients vanish unless $a+b+|\beta|$ is odd.  Since all reservoir degrees $d\in D$ are odd, $a+b+\sum_d d\beta_d$ has the same parity as $a+b+|\beta|$.  Thus only odd weighted degrees appear.
\end{proof}

\subsection{The third/fifth specialization}\label{subsec:third-fifth-coefficients}

\begin{definition}[Third/fifth coefficient data]\label{def:third-fifth-coeff-data}
For the scheme $D=\{3,5\}$, write
\[
  \alpha={s_3\over S},\qquad \beta={s_5\over S},
  \qquad S=(s_3^2+s_5^2)^{1/2}.
\]
With the numerical parameters \eqref{eq:params},
\[
  S=0.345068689589145\ldots,
  \quad
  \alpha=0.988241617650167\ldots,
  \quad
  \beta=0.152900311131734\ldots .
\]
\end{definition}

For $D=\{3,5\}$, Theorem~\ref{thm:coefficient-formula} becomes
\begin{equation}\label{eq:35-bm}
  b_m={\pi\over2}
  \sum_{\substack{a,b,\beta_3,\beta_5\in\Nzero\\ a+b+3\beta_3+5\beta_5=m}}
  (-1)^{b+\beta_3}
  {k!\over \beta_3!\beta_5!}
  \alpha^{2\beta_3}\beta^{2\beta_5}
  C_{a,b,k}^2,
  \qquad k=\beta_3+\beta_5.
\end{equation}
This is the formula used for the numerical head and tail certification in Section~\ref{sec:certification}.

\section{Finite-dimensional realization of the limiting certificate}\label{sec:transfer}

This section is Section~\ref{sec:transfer} in the proof architecture of Section~\ref{subsec:architecture}.

This section is the rigorous replacement for the informal statement that finite schemes approximate the limiting \(b\)-series.  Coefficientwise convergence is true, but it is not strong enough to transfer the inverse-majorant condition, which is an infinite nonlinear condition on the inverse series.  We prove a stronger analytic statement.

\subsection{Locally uniform convergence of finite correlations}

Let \(H_N\) denote the arcsine-normalized correlation function of the finite scheme \((f_N,g_N)\), with coordinatewise Gaussian correlation \(t\):
\[
  H_N(t)={\pi\over2}\E[f_N(\mathcal X)g_N(\mathcal Y_t)].
\]
Here \(\mathcal X,\mathcal Y_t\in\R^{2+|D|N}\) are standard Gaussian vectors with \(\E[\mathcal X_i\mathcal Y_{t,i}]=t\) coordinatewise.

\begin{proposition}[Local uniform convergence]\label{prop:loc-unif}
The functions \(H_N\) extend holomorphically to \(\D\), satisfy \(|H_N(z)|\le\pi/2\) on \(\D\), and converge locally uniformly on \(\D\) to the limiting weighted correlation \(H\).  Consequently, if
\[
  H_N(z)=\sum_{m\ge1}b_{m,N}z^m,
  \qquad
  H(z)=\sum_{m\ge1}b_mz^m,
\]
then \(b_{m,N}\to b_m\) for every fixed \(m\).
\end{proposition}

\begin{proof}
For fixed \(N\), expand \(f_N\) and \(g_N\) in the orthonormal tensor Hermite basis on \(\R^{2+|D|N}\):
\[
  f_N=\sum_\alpha \widehat f_N(\alpha)\mathcal H_\alpha,
  \qquad
  g_N=\sum_\alpha \widehat g_N(\alpha)\mathcal H_\alpha.
\]
Mehler's formula gives, for real \(t\in(-1,1)\),
\[
  {2\over\pi}H_N(t)=\sum_\alpha \widehat f_N(\alpha)\widehat g_N(\alpha)t^{|\alpha|}.
\]
The same series is absolutely convergent for \(|z|<1\), since Cauchy-Schwarz gives
\begin{align*}
  \sum_\alpha |\widehat f_N(\alpha)\widehat g_N(\alpha)| |z|^{|\alpha|}
  &\le
  \left(\sum_\alpha |\widehat f_N(\alpha)|^2 |z|^{|\alpha|}\right)^{1/2}
  \left(\sum_\alpha |\widehat g_N(\alpha)|^2 |z|^{|\alpha|}\right)^{1/2}\\
  &\le \|f_N\|_2\|g_N\|_2=1.
\end{align*}
Thus \(H_N\) is holomorphic on \(\D\) and \(|H_N(z)|\le \pi/2\).

We next identify the pointwise limit on the real interval.  Fix \(t\in(-1,1)\).  For every reservoir degree \(d\), the Hermite covariance identity gives
\[
  \E[H_d(U_{d,j})H_d(U_{d,j}')]=t^d.
\]
The multivariate central limit theorem yields
\[
  \bigl(P_{d,N},P_{d,N}'\bigr)_{d\in D}
  \Rightarrow
  \bigl(R_d,R_d'\bigr)_{d\in D},
\]
where \((R_d,R_d')\) is a Gaussian pair with correlation \(t^d\), and the pairs are independent over distinct \(d\).  This convergence is joint with \((Z,Z',X,X')\), which are independent of the reservoirs.

The limiting discontinuity set of the product of signs is contained in the union of the two threshold hypersurfaces
\begin{align*}
  z+\eta H_3(x)+\sum_{d\in D}s_dr_d&=0,\\
  z'-\eta H_3(x')+\sum_{d\in D}\sigma_ds_dr_d'&=0.
\end{align*}
This union has limiting probability zero: for example, after conditioning on all variables except \(Z\), the first threshold has probability zero because \(Z\) has a density, and similarly for the second threshold using \(Z'\).  The bounded continuous mapping theorem therefore gives \(H_N(t)\to H(t)\) for every real \(t\in(-1,1)\).

The family \(\{H_N\}\) is locally bounded on \(\D\) and converges pointwise on a set with an accumulation point.  Vitali's theorem gives locally uniform convergence on \(\D\) to the unique holomorphic function agreeing with the real limiting correlation.  Coefficient convergence follows from Cauchy's formula.
\end{proof}

\begin{remark}
Proposition~\ref{prop:loc-unif} justifies the limiting coefficient formula as the fixed-degree limit of genuine finite-dimensional schemes.  But it does not by itself imply that the finite inverse series satisfy \(\sum_n|A_{n,N}|\gamma^n\le1\).  Tail control for the inverse series requires the inverse-branch stability theorem below.
\end{remark}

\subsection{Analytic stability of inverse branches}

\begin{theorem}[Finite inverse-majorant transfer]\label{thm:finite-transfer}
Let \(F_N,F\in\Hol(\D)\), assume \(F_N\to F\) locally uniformly on \(\D\), and assume \(F_N(0)=F(0)=0\).  Suppose that \(G\) is a holomorphic inverse branch of \(F\) on a neighborhood of \(\ol{D(0,R_0)}\), with values in \(\D\):
\[
  F(G(\zeta))=\zeta,
  \qquad |\zeta|\le R_0.
\]
Write
\[
  G(\zeta)=\sum_{n\ge1}A_n\zeta^n.
\]
Let \(0<\gamma<R_0\), and assume the strict majorant inequality
\[
  \sum_{n\ge1}|A_n|\gamma^n<1.
\]
Then, for all sufficiently large \(N\), \(F_N\) has a holomorphic inverse branch at the origin,
\[
  G_N(\zeta)=F_N^{-1}(\zeta)=\sum_{n\ge1}A_{n,N}\zeta^n,
\]
defined on a neighborhood of \(\gamma\ol\D\), and
\[
  \sum_{n\ge1}|A_{n,N}|\gamma^n<1.
\]
Moreover \(A_{n,N}\to A_n\) for every fixed \(n\).
\end{theorem}

\begin{proof}
Choose nested radii
\[
  \gamma<\rho<r<R<R_0.
\]
Let
\[
  \Omega=G(D(0,R)).
\]
Since \(G\) is a conformal inverse branch on a neighborhood of \(\ol{D(0,R)}\), the boundary \(\partial\Omega\) is the image of \(|\zeta|=R\).  For \(t\in\partial\Omega\), we have \(|F(t)|=R\).  Therefore, for all \(|\zeta|\le r\),
\[
  |F(t)-\zeta|\ge R-r,
  \qquad t\in\partial\Omega.
\]
The compact contour \(\partial\Omega\) is contained in \(\D\).  By locally uniform convergence, for all large \(N\),
\[
  \sup_{t\in\partial\Omega}|F_N(t)-F(t)|<R-r.
\]
Rouche's theorem then implies that, for every \(|\zeta|\le r\), the equations \(F_N(t)=\zeta\) and \(F(t)=\zeta\) have the same number of solutions in \(\Omega\), counted with multiplicity.  The equation \(F(t)=\zeta\) has exactly one solution in \(\Omega\), namely \(t=G(\zeta)\).  Hence \(F_N(t)=\zeta\) also has exactly one solution in \(\Omega\).  In particular, since \(F_N(0)=0\), the zero at \(0\) is simple, and these solutions define a holomorphic inverse branch \(G_N\) on \(D(0,r)\), with values in \(\Omega\).

The inverse branches converge uniformly on \(\ol{D(0,r)}\).  Indeed, by the argument-principle formula,
\[
  G_N(\zeta)= {1\over 2\pi i}\int_{\partial\Omega}
  {tF_N'(t)\over F_N(t)-\zeta}\,dt,
  \qquad |\zeta|\le r.
\]
The denominators are uniformly bounded away from zero on \(\partial\Omega\times\ol{D(0,r)}\), and \(F_N'\to F'\) uniformly on \(\partial\Omega\).  Thus \(G_N\to G\) uniformly on \(\ol{D(0,r)}\), which implies fixed inverse-coefficient convergence.

It remains to control the tail.  Put
\[
  B=\sup_{t\in\ol\Omega}|t|<\infty.
\]
Since \(G_N(D(0,r))\subseteq\Omega\), Cauchy's estimate on the circle \(|\zeta|=\rho\) gives
\[
  |A_{n,N}|\le B\rho^{-n},
  \qquad n\ge1,
\]
for all large \(N\).  Let
\[
  \delta=1-\sum_{n\ge1}|A_n|\gamma^n>0.
\]
Choose \(M\) so large that
\[
  B\sum_{n>M}\left({\gamma\over\rho}\right)^n< {\delta\over3}.
\]
For all sufficiently large \(N\), fixed-coefficient convergence gives
\[
  \sum_{n\le M}|A_{n,N}|\gamma^n
  \le \sum_{n\le M}|A_n|\gamma^n+{\delta\over3}.
\]
Therefore
\begin{align*}
  \sum_{n\ge1}|A_{n,N}|\gamma^n
  &\le \sum_{n\le M}|A_n|\gamma^n+{\delta\over3}
  +B\sum_{n>M}\left({\gamma\over\rho}\right)^n\\
  &< \sum_{n\ge1}|A_n|\gamma^n+{2\delta\over3}<1.
\end{align*}
\end{proof}

\begin{corollary}[Finite schemes from a limiting certificate]\label{cor:finite-scheme}
Assume the limiting weighted correlation \(H\) has a holomorphic inverse branch
\[
  G(\zeta)=H^{-1}(\zeta)=\sum_{n\ge1}A_n\zeta^n
\]
with values in \(\D\) on a neighborhood of \(\ol{D(0,R_0)}\).  Let \(0<\gamma<R_0\), and suppose
\[
  \sum_{n\ge1}|A_n|\gamma^n<1.
\]
Then, for all sufficiently large reservoir lengths \(N\), the genuine finite-dimensional scheme \((f_N,g_N)\) satisfies the inverse-majorant condition at radius \(\gamma\).  Consequently
\[
  \KG\le {\pi\over2\gamma}.
\]
\end{corollary}

\begin{proof}
Apply Proposition~\ref{prop:loc-unif} and Theorem~\ref{thm:finite-transfer} with \(F_N=H_N\) and \(F=H\), then invoke Theorem~\ref{thm:krivine-criterion} for the finite scheme.
\end{proof}

\begin{corollary}[Endpoint and margin]\label{cor:margin}
If the limiting inverse branch is holomorphic with values in \(\D\) on a neighborhood of \(\Gamma\ol\D\) and
\[
  \sum_{n\ge1}|A_n|\Gamma^n\le1,
\]
then every \(\gamma<\Gamma\) transfers to all sufficiently high finite-dimensional schemes.  If the inequality at \(\Gamma\) is strict and the same analytic neighborhood is available, then \(\Gamma\) itself transfers.
\end{corollary}

\begin{proof}
For \(\gamma<\Gamma\), the majorant is strict because \(A_1\neq0\):
\[
  \sum_{n\ge1}|A_n|\gamma^n<\sum_{n\ge1}|A_n|\Gamma^n\le1.
\]
Apply Corollary~\ref{cor:finite-scheme}.  The strict endpoint case is the same argument with \(\gamma=\Gamma\).
\end{proof}

\section{A Wiener-algebra certificate for the limiting inverse}\label{sec:wiener}

This section is Section~\ref{sec:wiener} in the proof architecture of Section~\ref{subsec:architecture}.

We now turn to the limiting inverse-majorant condition.  This is where the computation enters.  The key point is that the nonlinear part of \(H\) is very small in the Wiener algebra.


\subsection{The Wiener algebra}\label{subsec:wiener-algebra}

\begin{definition}[Wiener-algebra data]\label{def:wiener-data}
Let $\calA$ be the Wiener algebra of absolutely convergent power series on the unit disk:
\[
  U(w)=\sum_{n\ge0}u_nw^n,
  \qquad
  \|U\|_1=\sum_{n\ge0}|u_n|<\infty.
\]
Then $\|UV\|_1\le\|U\|_1\|V\|_1$.  Write
\begin{equation}\label{eq:H-linear-nonlinear}
  H(t)=b_1t+E(t),
  \qquad
  E(t)=\sum_{m\ge3}b_mt^m,
\end{equation}
and define
\begin{equation}\label{eq:eta-L}
  \eta_0=\sum_{m\ge3}|b_m|,
  \qquad
  L=\sum_{m\ge3}m|b_m|.
\end{equation}
\end{definition}

On the unit ball of $\calA$, the nonlinear part is $L$-Lipschitz:
\begin{equation}\label{eq:E-lipschitz}
  \|E(U)-E(V)\|_1\le L\|U-V\|_1,
  \qquad \|U\|_1,\|V\|_1\le1.
\end{equation}
This follows from $\|U^m-V^m\|_1\le m\|U-V\|_1$ on the unit ball.

\begin{theorem}[A simple linear-dominance certificate]\label{thm:simple-wiener}
Assume
\[
  L<b_1,
  \qquad
  \gamma+\eta_0<b_1.
\]
Then the limiting inverse satisfies
\[
  \sum_{n\ge1}|A_n|\gamma^n<1,
  \qquad H^{-1}(\zeta)=\sum_{n\ge1}A_n\zeta^n.
\]
Moreover the inverse branch is holomorphic on a neighborhood of \(\gamma\ol\D\), and maps \(\gamma\D\) into \(\D\).
\end{theorem}

\begin{proof}
Seek a solution \(U(w)=H^{-1}(\gamma w)\) in the closed unit ball of \(\calA\) with \(U(0)=0\).  The inverse equation is equivalent to
\[
  U=\Phi(U):={\gamma w-E(U)\over b_1}.
\]
If \(\|U\|_1\le1\), then \(\|E(U)\|_1\le\eta_0\), hence
\[
  \|\Phi(U)\|_1\le {\gamma+\eta_0\over b_1}<1.
\]
By \eqref{eq:E-lipschitz}, \(\Phi\) is a contraction with constant \(L/b_1<1\).  Banach's fixed point theorem gives a unique fixed point \(U\in\calA\), and
\[
  \|U\|_1=\sum_{n\ge1}|A_n|\gamma^n<1.
\]
Since \(H(U(w))=\gamma w\) and \(U'(0)=\gamma/b_1\neq0\), this fixed point agrees near zero with the local inverse branch.  It remains to note that the branch has a little analytic room around the closed disk.  Put \(q=\|U\|_1<1\), and choose \(q<q_*<1\).  For \(u,v\in q_*\ol\D\),
\[
  |H(u)-H(v)|\ge b_1|u-v|-\sum_{m\ge3}|b_m||u^m-v^m|
  \ge (b_1-L)|u-v|.
\]
Thus \(H\) is injective on \(q_*\D\), and \(H'\) has no zero there.  Since \(U(\ol\D)\subset q\ol\D\subset q_*\D\), the set \(\gamma\ol\D=H(U(\ol\D))\) is contained in the open set \(H(q_*\D)\), on which the inverse branch is holomorphic.
\end{proof}

The simple criterion is conceptually useful but leaves numerical slack.  The sharper certificate uses a short inverse head plus the same Lipschitz estimate.


\subsection{A posteriori inverse-tail certificate}\label{subsec:aposteriori}

\begin{definition}[Inverse head and residual]\label{def:inverse-head-residual}
Let $G(\zeta)=H^{-1}(\zeta)$, and set $U(w)=G(\gamma w)$.  For an odd cutoff $N$, define the inverse head
\[
  P_N(w)=\sum_{1\le n\le N}A_n\gamma^n w^n,
  \qquad
  p_N=\|P_N\|_1.
\]
For a coefficient cutoff $K$, define
\[
  H_{\le K}(t)=\sum_{m\le K}b_mt^m,
  \qquad
  \eps_K=\sum_{m>K}|b_m|.
\]
Assuming $K\ge N$, the coefficients of $H_{\le K}(P_N)-\gamma w$ through degree $N$ vanish by construction of the inverse head.  The residual tail is
\begin{equation}\label{eq:rNK}
  r_{N,K}=\left\|T_{>N}\bigl(H_{\le K}(P_N(w))-\gamma w\bigr)\right\|_1,
\end{equation}
where $T_{>N}$ discards coefficients of degree at most $N$.
\end{definition}

\begin{theorem}[A posteriori inverse-tail certificate]\label{thm:aposteriori}
Assume \(L<b_1\), \(p_N<1\), and define
\begin{equation}\label{eq:RN}
  R_N={r_{N,K}+\eps_K\over b_1-L}.
\end{equation}
If
\[
  p_N+R_N<1,
\]
then the limiting inverse satisfies
\[
  \sum_{n\ge1}|A_n|\gamma^n<1.
\]
Moreover \(H^{-1}\) has an inverse branch on a neighborhood of \(\gamma\ol\D\), and its values on \(\gamma\D\) lie in \(\D\).
\end{theorem}

\begin{proof}
We solve the inverse equation in the form \(U=P_N+Q\).  Since
\[
  H(P_N+Q)=b_1(P_N+Q)+E(P_N+Q),
\]
the equation \(H(P_N+Q)=\gamma w\) is equivalent to
\begin{equation}\label{eq:Q-fixed}
  Q={\gamma w-H(P_N)-E(P_N+Q)+E(P_N)\over b_1}.
\end{equation}
Because \(p_N<1\), we have
\[
  \|H(P_N)-H_{\le K}(P_N)\|_1\le \eps_K.
\]
Also, by the definition of \(P_N\), the coefficients of \(H_{\le K}(P_N)-\gamma w\) up to degree \(N\) vanish; hence
\[
  \|H(P_N)-\gamma w\|_1\le r_{N,K}+\eps_K.
\]
Consider the right side of \eqref{eq:Q-fixed} on the ball \(\|Q\|_1\le R_N\).  If \(p_N+R_N<1\), then \(\|P_N+Q\|_1\le1\), so by \eqref{eq:E-lipschitz},
\[
  \|E(P_N+Q)-E(P_N)\|_1\le L\|Q\|_1.
\]
Thus the right side of \eqref{eq:Q-fixed} has norm at most
\[
  {r_{N,K}+\eps_K+LR_N\over b_1}=R_N.
\]
The same Lipschitz estimate gives contraction constant \(L/b_1<1\).  Hence there is a unique solution \(Q\) in this ball.  Therefore \(U=P_N+Q\) satisfies \(H(U)=\gamma w\) and
\[
  \|U\|_1\le p_N+R_N<1.
\]
As before, this analytic solution agrees near zero with the local inverse branch, proving the claimed inverse-majorant inequality.  Finally set \(q=p_N+R_N<1\) and choose \(q<q_*<1\).  The same estimate
\[
  |H(u)-H(v)|\ge (b_1-L)|u-v|,
  \qquad u,v\in q_*\ol\D,
\]
shows that \(H\) is injective and locally biholomorphic on \(q_*\D\).  Since \(U(\ol\D)\subset q\ol\D\), the disk \(\gamma\ol\D=H(U(\ol\D))\) is contained in the open set \(H(q_*\D)\).  Hence the inverse branch is holomorphic on a neighborhood of \(\gamma\ol\D\).
\end{proof}


\section{Certifying the coefficient head and tail}\label{sec:certification}

This section is the final proof section in the roadmap of Section~\ref{subsec:architecture}.  It records the coefficient inequalities needed for the limiting inverse certificate, proves Theorem~\ref{thm:main} from those inequalities, and then explains how the remaining head and tail computations should be certified.  The interval-arithmetic details will be inserted here in the final version.

\subsection{Certificate module and proof of the main theorem}\label{subsec:certificate}

The following is the only part of the draft that remains to be replaced by a formal interval-arithmetic proof.  The quantities are all finite or explicitly tailed through the coefficient formula \eqref{eq:35-bm}.

\begin{assumption}[Coefficient and inverse-head certificate]\label{ass:certificate}
For the parameters \eqref{eq:params}, the limiting coefficient sequence \(H(t)=\sum b_mt^m\) satisfies
\begin{align}
  b_1&>0.88157380,\label{eq:cert-b1}\\
  \eps_{41}:=\sum_{m>41}|b_m|&<10^{-8},\label{eq:cert-eps41}\\
  L:=\sum_{m\ge3}m|b_m|&<1.58\cdot10^{-4}.\label{eq:cert-L}
\end{align}
Moreover, for \(\gammazero=0.8815624\), \(N=21\), and \(K=41\), the inverse head and residual defined above satisfy
\begin{align}
  p_{21}&<0.9999992590,\label{eq:cert-p21}\\
  r_{21,41}&<5.574\cdot10^{-7}.\label{eq:cert-r2141}
\end{align}
\end{assumption}

\begin{proposition}[Limiting admissibility from the certificate]\label{prop:lim-admissible}
Assuming Certificate Module~\ref{ass:certificate}, the limiting third/fifth model is admissible at \(\gammazero=0.8815624\):
\[
  \sum_{n\ge1}|A_n|\gammazero^n<1,
  \qquad H^{-1}(\zeta)=\sum_{n\ge1}A_n\zeta^n.
\]
The inverse branch is holomorphic on a neighborhood of \(\gammazero\ol\D\), and maps \(\gammazero\D\) into \(\D\).
\end{proposition}

\begin{proof}
Apply Theorem~\ref{thm:aposteriori} with \(N=21\), \(K=41\), and \(\gamma=\gammazero\).  By the certificate,
\[
  R_{21}\le {5.574\cdot10^{-7}+10^{-8}\over 0.88157380-1.58\cdot10^{-4}}
  <6.45\cdot10^{-7}.
\]
Therefore
\[
  p_{21}+R_{21}<0.9999992590+0.000000645<1.
\]
The conclusion follows.
\end{proof}

\begin{proof}[Proof of Theorem~\ref{thm:main}]
By Proposition~\ref{prop:lim-admissible}, the limiting model satisfies a strict inverse-majorant certificate at \(\gammazero\), with an inverse branch holomorphic on a neighborhood of \(\gammazero\ol\D\) and with values in \(\D\) on \(\gammazero\D\).  Corollary~\ref{cor:finite-scheme} transfers this certificate to all sufficiently high finite-dimensional third/fifth-chaos schemes \eqref{eq:finite-f-35}--\eqref{eq:finite-g-35}.  Applying the Krivine criterion, Theorem~\ref{thm:krivine-criterion}, to such a finite scheme gives
\[
  \KG\le {\pi\over2\gammazero}=1.7818322637\ldots .
\]
Since
\[
  {\pi\over 2\log(1+\sqrt2)}-{\pi\over2\gammazero}
  =3.817144799\cdot10^{-4}\ldots,
\]
the advertised hyperplane-improvement form follows.
\end{proof}


\subsection{Numerical head of the limiting model}\label{subsec:numerical-head}

For the parameters \eqref{eq:params}, one has
\[
  S=0.345068689589145\ldots,
  \quad
  \alpha=0.988241617650167\ldots,
  \quad
  \beta=0.152900311131734\ldots .
\]
The following table gives the beginning of the limiting coefficient sequence.  These are floating-point values included to explain the numerics; the rigorous proof uses interval enclosures in Certificate Module~\ref{ass:certificate}.

\begin{center}
\small
\begin{tabular}{r r@{\qquad}r r}
\toprule
\(m\) & \multicolumn{1}{c}{\(b_m\)} & \(m\) & \multicolumn{1}{c}{\(b_m\)}\\
\midrule
1  & \( 8.815738220496\cdot10^{-1}\) & 23 & \( 3.133017408224\cdot10^{-7}\)\\
3  & \( 3.453873298476\cdot10^{-9}\) & 25 & \( 1.468544711235\cdot10^{-7}\)\\
5  & \(-9.043545773357\cdot10^{-11}\) & 27 & \( 6.021901119747\cdot10^{-8}\)\\
7  & \( 1.171432343870\cdot10^{-7}\) & 29 & \( 2.058498871742\cdot10^{-8}\)\\
9  & \(-1.494586697161\cdot10^{-9}\) & 31 & \( 3.428828624581\cdot10^{-9}\)\\
11 & \(-5.743183500533\cdot10^{-6}\) & 33 & \(-2.837596359323\cdot10^{-9}\)\\
13 & \(-1.448951029677\cdot10^{-6}\) & 35 & \(-3.866295291032\cdot10^{-9}\)\\
15 & \( 9.441806267864\cdot10^{-7}\) & 37 & \(-2.977047658664\cdot10^{-9}\)\\
17 & \( 1.104203285670\cdot10^{-6}\) & 39 & \(-1.857241836324\cdot10^{-9}\)\\
19 & \( 8.451950461437\cdot10^{-7}\) & 41 & \(-1.033472390529\cdot10^{-9}\)\\
21 & \( 5.630558521346\cdot10^{-7}\) & & \\
\bottomrule
\end{tabular}
\end{center}

The cancellation of \(b_3\) and \(b_5\) is the visible reason the scheme beats the hyperplane radius.  If
\[
  H(t)=b_1t+b_3t^3+b_5t^5+\cdots,
  \qquad
  H^{-1}(\zeta)=A_1\zeta+A_3\zeta^3+A_5\zeta^5+\cdots,
\]
then formal series reversion gives
\begin{equation}\label{eq:low-inverse}
  A_1={1\over b_1},
  \qquad
  A_3=-{b_3\over b_1^4},
  \qquad
  A_5={3b_3^2-b_1b_5\over b_1^7}.
\end{equation}
Thus killing \(b_3\) and \(b_5\) kills the first two nonlinear inverse coefficients.

For reference, the inverse coefficients begin as follows:
\begin{center}
\small
\begin{tabular}{r r@{\qquad}r r}
\toprule
\(m\) & \multicolumn{1}{c}{\(A_m\)} & \(m\) & \multicolumn{1}{c}{\(A_m\)}\\
\midrule
1  & \( 1.134334952999\) & 13 & \( 8.461201717465\cdot10^{-6}\)\\
3  & \(-5.718362197765\cdot10^{-9}\) & 15 & \(-7.094407911980\cdot10^{-6}\)\\
5  & \( 1.926579003756\cdot10^{-10}\) & 17 & \(-1.067574204152\cdot10^{-5}\)\\
7  & \(-3.211054870543\cdot10^{-7}\) & 19 & \(-1.051441386906\cdot10^{-5}\)\\
9  & \( 5.271516751153\cdot10^{-9}\) & 21 & \(-9.006179634605\cdot10^{-6}\)\\
11 & \( 2.606444408195\cdot10^{-5}\) & & \\
\bottomrule
\end{tabular}
\end{center}


\subsection{What must be certified}

The proof above requires only the inequalities in Certificate Module~\ref{ass:certificate}.  A convenient implementation can be organized as follows.

\begin{enumerate}[label=\textup{C\arabic*.}, leftmargin=2.7em]
\item Certify the low coefficients through degree \(41\), especially
\[
  b_1>0.88157380,
  \qquad
  \sum_{3\le m\le41}|b_m|<1.1328\cdot10^{-5},
  \qquad
  \sum_{3\le m\le41}m|b_m|<1.574\cdot10^{-4}.
\]

\item Certify the unweighted coefficient tail
\[
  \sum_{m>41}|b_m|<10^{-8}.
\]

\item Certify the weighted nonlinear mass
\[
  \sum_{m\ge3}m|b_m|<1.58\cdot10^{-4}.
\]

\item Using the enclosed coefficients through degree \(41\), compute the inverse coefficients through degree \(21\) and certify
\[
  p_{21}=\sum_{n\le21}|A_n|\gammazero^n<0.9999992590.
\]
Then compute the residual polynomial
\[
  T_{>21}\bigl(H_{\le41}(P_{21}(w))-\gammazero w\bigr)
\]
and certify
\[
  r_{21,41}<5.574\cdot10^{-7}.
\]
\end{enumerate}

Items C1 and C4 are finite arithmetic once the coefficients are enclosed.  The true analytic tail work is C2 and C3.

\subsection{Why the original coefficient tail is easier than the inverse tail}

The inverse-tail asks for coefficients of \(H^{-1}\), whose branch approaches the unit boundary in the intended application.  A naive Cauchy estimate for the inverse therefore needs many coefficients and a delicate complex-domain bound.  The coefficient tail of \(H\) is more structured.  Formula \eqref{eq:35-bm} expresses every \(b_m\) as a finite weighted-degree sum of explicit real two-dimensional integrals \eqref{eq:C-abk}.  Thus the tail can be attacked directly by Gaussian quadrature with certified remainder, Taylor models on boxes plus Gaussian tails, or Sobolev/Hermite tail estimates for the functions appearing in \eqref{eq:C-abk}.

The observed signed coefficients beyond degree \(41\) are already tiny in floating point:
\[
  \sum_{43\le m\le201}|b_m|\approx 9.93\cdot10^{-10},
  \qquad
  \sum_{43\le m\le201}m|b_m|\approx 5.22\cdot10^{-8}.
\]
A practical route is therefore to enclose all coefficients through degree \(201\), verify the finite sums, and then use a general analytic Hermite-tail estimate for \(m>201\).

\section{Discussion and future directions}\label{sec:discussion}

This discussion section is Section~\ref{sec:discussion} in the proof architecture of Section~\ref{subsec:architecture}.  It explains the parameter choices, nearby extensions, and finite-dimensional effectiveness questions.

\subsection{Why degrees three and five}

The low-order inverse relations \eqref{eq:low-inverse} explain the parameter search.  The hyperplane scheme has \(H(t)=\arcsin t\), so its inverse \(\sin\zeta\) has alternating nonlinear coefficients.  To increase the admissible radius, the first objective is to suppress the cubic inverse obstruction.  An anti-aligned third-chaos reservoir changes \(b_3\) in the correct direction.  Once \(b_3\) is controlled, the next obstruction is quintic, and an aligned fifth-chaos reservoir gives the next degree of freedom.

The small finite perturbation \(\eta H_3(X)\) in the base threshold is also important.  A purely reservoir-Gaussian threshold is too close to a hyperplane after rotation; the \(X\)-dependent cubic term changes the two-dimensional integrals \(C_{a,b,k}\) and gives the higher coefficients the right shape.

\subsection{A degree \texorpdfstring{\(7/9\)}{7/9} extension}

The coefficient formula immediately extends to additional odd reservoirs.  For instance, taking
\[
  D=\{3,5,7,9\},
  \qquad
  \sigma_d=(-1)^{(d-1)/2},
\]
keeps the same formula \eqref{eq:general-bm}.  A preliminary numerical search gives
\begin{align*}
  \eta&=0.12569605,&
  s_3&=0.34199778,&
  s_5&=0.05153822,\\
  s_7&=0.00698313,&
  s_9&=0.00090252.
\end{align*}
For the corresponding limiting model, the estimated value is
\[
  \gamma\approx0.88157525281,
  \qquad
  \KG\lesssim1.78180628572.
\]
This extension is not used in the main theorem.  It is included to indicate that the third/fifth scheme is not isolated: the same analytic and finite-transfer machinery applies to a broader weighted-chaos family.

\subsection{Effective dimensions}

Corollary~\ref{cor:finite-scheme} is qualitative: it proves the existence of sufficiently large finite reservoir length \(N\), but it does not give a useful dimension bound.  For the Grothendieck constant as a mathematical constant, qualitative existence of a finite scheme is enough.  For an explicit algorithmic rounding implementation, one should quantify the convergence \(H_N\to H\) on the contour used in Theorem~\ref{thm:finite-transfer}.

The natural route is a collision expansion.  For fixed reservoir degree \(d\) and fixed \(k\), the polynomial \(H_k(P_{d,N})\) has a leading distinct-index component
\[
  \sqrt{k!}\,N^{-k/2}
  \sum_{1\le i_1<\cdots<i_k\le N}
  H_d(U_{d,i_1})\cdots H_d(U_{d,i_k}),
\]
which lies in pure Gaussian chaos degree \(dk\) and has \(L^2\)-norm squared \((N)_k/N^k\to1\).  The remaining collision terms have vanishing \(L^2\)-mass for fixed \(d,k\).  This explains at the coefficient level why a limiting reservoir Hermite mode of order \(k\) contributes correlation power \(t^{dk}\).  Making this expansion uniform to the required degree would produce explicit finite-dimensional estimates.


\appendix

\section{Krivine preprocessing and mixed rounding}\label{app:preprocessing}

This appendix contains the standard preprocessing and projection arguments behind Krivine rounding schemes.  The argument is the rounding construction of Krivine and BMMN~\cite{Krivine1977,BMMN2011}, written in the variance-one Gaussian normalization used throughout this paper.  It is included here to keep the body of the paper focused on the new chaos construction.

For a real matrix $A=(a_{ij})\in\R^{m\times n}$, write
\[
  \SDP(A):=\sup_{\|x_i\|=\|y_j\|=1}
  \left|\sum_{i=1}^m\sum_{j=1}^n a_{ij}\ip{x_i}{y_j}\right|
\]
and
\[
  \OPT(A):=\max_{\eps_i,\delta_j\in\{\pm1\}}
  \left|\sum_{i=1}^m\sum_{j=1}^n a_{ij}\eps_i\delta_j\right|.
\]
Thus proving $\SDP(A)\le K\OPT(A)$ for all $A$ proves $\KG\le K$.

\begin{theorem}[Krivine preprocessing and projection]\label{thm:krivine-criterion}
Let $f,g:\R^k\to\{-1,1\}$ be odd measurable functions, and let $H=H_{f,g}$ be the arcsine-normalized correlation function
\begin{equation}\label{eq:pure-correlation}
  H(t)={\pi\over2}\E\left[f(G_1)
  g\left(tG_1+\sqrt{1-t^2}G_2\right)\right],
  \qquad -1<t<1,
\end{equation}
where $G_1,G_2$ are independent standard Gaussian vectors in $\R^k$.  Suppose that $H$ has a local inverse at the origin,
\[
  H^{-1}(\zeta)=\sum_{r\ge1}a_r\zeta^r,
\]
which is holomorphic on a neighborhood of $\gamma\ol\D$, and suppose that
\[
  \sum_{r\ge1}|a_r|\gamma^r\le1.
\]
Then
\[
  \KG\le {\pi\over2\gamma}.
\]
\end{theorem}

\begin{proof}
It is enough to prove the conclusion with any $0<\gamma'<\gamma$ in place of $\gamma$, and then let $\gamma'\uparrow\gamma$.  We therefore fix such a $\gamma'$.  Put
\[
  M_{\gamma'}:=\sum_{r\ge1}|a_r|(\gamma')^r<1.
\]
Let $A=(a_{ij})\in\R^{m\times n}$.  Choose unit vectors $x_1,\ldots,x_m,y_1,\ldots,y_n$ such that
\[
  \sum_{i,j}a_{ij}\ip{x_i}{y_j}
\]
is arbitrarily close to $\SDP(A)$; if necessary, replace all $y_j$ by $-y_j$ so that this quantity is nonnegative.  We suppress this arbitrarily small error, since it disappears at the end.

Let
\[
  \calK:=\left(\bigoplus_{r\ge1}\calH^{\otimes r}\right)\oplus \R e_L\oplus \R e_R,
\]
where $\calH$ is the Hilbert space containing the original vectors, and where $e_L,e_R$ are unit vectors orthogonal to each other and to all tensor summands.  Define
\[
  I(x):=\bigoplus_{r\ge1}\sqrt{|a_r|}(\gamma')^{r/2}x^{\otimes r}
  \oplus \sqrt{1-M_{\gamma'}}e_L\oplus0
\]
and
\[
  J(y):=\bigoplus_{r\ge1}\sgn(a_r)\sqrt{|a_r|}(\gamma')^{r/2}y^{\otimes r}
  \oplus0\oplus \sqrt{1-M_{\gamma'}}e_R,
\]
with the convention that the summand corresponding to $a_r=0$ is zero.  Then $I(x)$ and $J(y)$ are unit vectors whenever $x,y$ are unit vectors.  Moreover,
\[
  \ip{I(x)}{J(y)}=\sum_{r\ge1}a_r(\gamma')^r\ip{x}{y}^r
  =H^{-1}(\gamma'\ip{x}{y}).
\]
The first equality uses
\[
  \ip{x^{\otimes r}}{y^{\otimes r}}=\ip{x}{y}^r,
\]
and the second equality uses the Taylor expansion of the inverse on $\gamma'\D$.

Set
\[
  u_i:=I(x_i),\qquad v_j:=J(y_j).
\]
The finitely many vectors $u_i,v_j$ span a finite-dimensional subspace of $\calK$.  Identify this subspace with $\R^N$.  Let $\Gamma:\R^N\to\R^k$ be a random Gaussian linear map, i.e. a $k\times N$ matrix whose entries are independent standard Gaussian random variables.  Define random signs
\[
  \eps_i:=f(\Gamma u_i),\qquad \delta_j:=g(\Gamma v_j).
\]
For each pair $(i,j)$, the vectors $\Gamma u_i$ and $\Gamma v_j$ are standard Gaussian vectors in $\R^k$ with coordinatewise correlation $\ip{u_i}{v_j}$.  Hence, by the definition of $H$,
\[
  \E[\eps_i\delta_j]={2\over\pi}H(\ip{u_i}{v_j}).
\]
Using the preprocessing identity,
\[
  H(\ip{u_i}{v_j})=H\left(H^{-1}(\gamma'\ip{x_i}{y_j})\right)=\gamma'\ip{x_i}{y_j}.
\]
Therefore
\[
  \E\left[\sum_{i,j}a_{ij}\eps_i\delta_j\right]
  ={2\over\pi}\sum_{i,j}a_{ij}H(\ip{u_i}{v_j})
  ={2\gamma'\over\pi}\sum_{i,j}a_{ij}\ip{x_i}{y_j}.
\]
Since the maximum over deterministic signs is at least the expectation of the randomized signs, we get
\[
  \OPT(A)\ge {2\gamma'\over\pi}\SDP(A).
\]
Thus
\[
  \SDP(A)\le {\pi\over2\gamma'}\OPT(A).
\]
Letting $\gamma'\uparrow\gamma$ gives
\[
  \SDP(A)\le {\pi\over2\gamma}\OPT(A).
\]
Since $A$ was arbitrary, $\KG\le \pi/(2\gamma)$.
\end{proof}

\begin{corollary}[Classical Krivine hyperplane bound]\label{cor:hyperplane-bound}
Let $h:\R\to\{-1,1\}$ be $h(x)=\sgn(x)$.  Then the hyperplane scheme $(h,h)$ gives
\[
  \KG\le {\pi\over 2\log(1+\sqrt2)}.
\]
\end{corollary}

\begin{proof}
For the hyperplane scheme, Grothendieck's identity gives
\[
  \E[\sgn(X)\sgn(Y)]={2\over\pi}\arcsin(t)
\]
whenever $X,Y$ are standard Gaussians with correlation $t$.  In our normalization,
\[
  H_{h,h}(t)=\arcsin(t).
\]
Thus
\[
  H_{h,h}^{-1}(\zeta)=\sin\zeta=\sum_{q\ge0}{(-1)^q\over(2q+1)!}\zeta^{2q+1}.
\]
Let
\[
  \rho_*:=\arsinh(1)=\log(1+\sqrt2).
\]
Then
\[
  \sum_{q\ge0}{\rho_*^{2q+1}\over(2q+1)!}=\sinh(\rho_*)=1.
\]
Theorem~\ref{thm:krivine-criterion} therefore applies with $\gamma=\rho_*$, giving
\[
  \KG\le {\pi\over2\rho_*}={\pi\over2\log(1+\sqrt2)}.
\]
\end{proof}

\begin{theorem}[Mixed Krivine preprocessing and projection]\label{thm:mixed-krivine}
Let $L\in\N$.  For each $\ell\in\{1,\ldots,L\}$, let $f_\ell,g_\ell:\R^{k_\ell}\to\{-1,1\}$ be odd measurable functions, and let $H_\ell$ be the corresponding arcsine-normalized correlation function:
\[
  H_\ell(t)={\pi\over2}\E\left[f_\ell(G_1)g_\ell\left(tG_1+\sqrt{1-t^2}G_2\right)\right].
\]
Let $\lambda_1,\ldots,\lambda_L\ge0$ satisfy $\sum_{\ell=1}^L\lambda_\ell=1$, and define the mixed function
\[
  H_\lambda(t):=\sum_{\ell=1}^L\lambda_\ell H_\ell(t).
\]
Suppose that $H_\lambda$ has a local inverse at the origin,
\[
  H_\lambda^{-1}(\zeta)=\sum_{r\ge1}a_r\zeta^r,
\]
which is holomorphic on a neighborhood of $\gamma\ol\D$, and suppose that
\[
  M_\lambda(\gamma):=\sum_{r\ge1}|a_r|\gamma^r\le1.
\]
Then
\[
  \KG\le {\pi\over2\gamma}.
\]
\end{theorem}

\begin{proof}
The preprocessing is identical to the proof of Theorem~\ref{thm:krivine-criterion}, but using the inverse coefficients of $H_\lambda^{-1}$.  As before, first fix $0<\gamma'<\gamma$, and set
\[
  M_{\lambda,\gamma'}:=\sum_{r\ge1}|a_r|(\gamma')^r<1.
\]
For unit vectors $x,y$, define
\[
  I(x):=\bigoplus_{r\ge1}\sqrt{|a_r|}(\gamma')^{r/2}x^{\otimes r}
  \oplus \sqrt{1-M_{\lambda,\gamma'}}e_L\oplus0
\]
and
\[
  J(y):=\bigoplus_{r\ge1}\sgn(a_r)\sqrt{|a_r|}(\gamma')^{r/2}y^{\otimes r}
  \oplus0\oplus\sqrt{1-M_{\lambda,\gamma'}}e_R.
\]
Then $I(x)$ and $J(y)$ are unit vectors and
\[
  \ip{I(x)}{J(y)}=H_\lambda^{-1}(\gamma'\ip{x}{y}).
\]
Now let $A=(a_{ij})$ be arbitrary, and choose unit vectors $x_i,y_j$ that nearly attain $\SDP(A)$, oriented so that $\sum_{i,j}a_{ij}\ip{x_i}{y_j}\ge0$.  Put
\[
  u_i:=I(x_i),\qquad v_j:=J(y_j).
\]
Identify the span of the finitely many $u_i,v_j$ with $\R^N$.  For each $\ell$, independently sample a Gaussian linear map
\[
  \Gamma_\ell:\R^N\to\R^{k_\ell}.
\]
Also sample a single global index $\Lambda\in\{1,\ldots,L\}$, independent of the Gaussian maps, with
\[
  \Prob[\Lambda=\ell]=\lambda_\ell.
\]
The word ``global'' is important: the same index $\Lambda$ is used for all $i,j$.  Define
\[
  \eps_i:=f_\Lambda(\Gamma_\Lambda u_i),
  \qquad
  \delta_j:=g_\Lambda(\Gamma_\Lambda v_j).
\]
Conditioning on $\Lambda=\ell$, the same rotation-invariance calculation as in the pure case gives
\[
  \E[\eps_i\delta_j\mid \Lambda=\ell]={2\over\pi}H_\ell(\ip{u_i}{v_j}).
\]
Averaging over $\Lambda$,
\[
  \E[\eps_i\delta_j]={2\over\pi}\sum_{\ell=1}^L\lambda_\ell H_\ell(\ip{u_i}{v_j})
  ={2\over\pi}H_\lambda(\ip{u_i}{v_j}).
\]
Since
\[
  \ip{u_i}{v_j}=H_\lambda^{-1}(\gamma'\ip{x_i}{y_j}),
\]
we obtain
\[
  \E[\eps_i\delta_j]={2\gamma'\over\pi}\ip{x_i}{y_j}.
\]
Therefore
\[
  \OPT(A)\ge
  \E\left[\sum_{i,j}a_{ij}\eps_i\delta_j\right]
  ={2\gamma'\over\pi}\sum_{i,j}a_{ij}\ip{x_i}{y_j}.
\]
Taking the supremum over the original vectors gives
\[
  \OPT(A)\ge {2\gamma'\over\pi}\SDP(A),
\]
and hence
\[
  \SDP(A)\le {\pi\over2\gamma'}\OPT(A).
\]
Letting $\gamma'\uparrow\gamma$ proves
\[
  \SDP(A)\le {\pi\over2\gamma}\OPT(A).
\]
Since $A$ was arbitrary, $\KG\le \pi/(2\gamma)$.
\end{proof}

\begin{remark}
Theorem~\ref{thm:krivine-criterion} is the special case $L=1$ of Theorem~\ref{thm:mixed-krivine}.  We separated the two statements because the pure case is the usual Krivine scheme, while the mixed case emphasizes the randomized choice of a common scheme $\Lambda$ before projection.
\end{remark}


\section*{Acknowledgments}
This draft is a proof-architecture document.  The final version should replace Certificate Module~\ref{ass:certificate} with a reproducible interval-arithmetic certificate and should record the precise code, parameter intervals, and rounding modes used in the computation.

\begin{thebibliography}{99}

\bibitem{AlonNaor2006}
N.~Alon and A.~Naor.
Approximating the cut-norm via Grothendieck's inequality.
\emph{SIAM Journal on Computing}, 35(4):787--803, 2006.

\bibitem{BMMN2011}
M.~Braverman, K.~Makarychev, Y.~Makarychev, and A.~Naor.
The Grothendieck constant is strictly smaller than Krivine's bound.
\emph{Forum of Mathematics, Pi}, 1:e4, 2013.
Conference version in FOCS 2011.  arXiv:1103.6161.


\bibitem{Davie1984}
A.~M. Davie.
Matrix norms related to Grothendieck's inequality.
In \emph{Banach spaces (Columbia, Mo., 1984)}, volume 1166 of Lecture Notes in Mathematics, pages 22--26. Springer, 1985.

\bibitem{Grothendieck1953}
A.~Grothendieck.
Resume de la theorie metrique des produits tensoriels topologiques.
\emph{Boletim da Sociedade de Matematica de Sao Paulo}, 8:1--79, 1953.

\bibitem{Johansson2017}
F.~Johansson.
Arb: efficient arbitrary-precision midpoint-radius interval arithmetic.
\emph{IEEE Transactions on Computers}, 66(8):1281--1292, 2017.


\bibitem{KhotNaor2011}
S.~Khot and A.~Naor.
Grothendieck-type inequalities in combinatorial optimization.
\emph{Communications on Pure and Applied Mathematics}, 65(7):992--1035, 2012.

\bibitem{LindenstraussPelczynski1968}
J.~Lindenstrauss and A.~Pe\l czy\'nski.
Absolutely summing operators in $L_p$-spaces and their applications.
\emph{Studia Mathematica}, 29:275--326, 1968.

\bibitem{LindenstraussTzafriri2013}
J.~Lindenstrauss and L.~Tzafriri.
\emph{Classical Banach Spaces I: Sequence Spaces}.
Springer, reprint edition, 2013.

\bibitem{Krivine1977}
J.-L.~Krivine.
Sur la constante de Grothendieck.
\emph{Comptes Rendus de l'Academie des Sciences de Paris, Serie A-B}, 284(8):A445--A446, 1977.

\bibitem{LiSahaXueEtAl2026}
A.~Li, R.~Saha, A.~Xue, A.~Klivans, P.~K. Kothari, R.~Meka, and S.~Chaudhuri.
The Grothendieck constant is less than \(\pi/(2\log(1+\sqrt2))-10^{-5}\).
arXiv:2606.03991, 2026.

\bibitem{NaorRegev2014}
A.~Naor and O.~Regev.
Krivine schemes are optimal.
\emph{Proceedings of the American Mathematical Society}, 142(12):4315--4320, 2014.


\bibitem{Reeds1991}
J.~A. Reeds.
A new lower bound on the real Grothendieck constant.
Unpublished manuscript, 1991.

\bibitem{Pisier2012}
G.~Pisier.
Grothendieck's theorem, past and present.
\emph{Bulletin of the American Mathematical Society}, 49(2):237--323, 2012.

\end{thebibliography}

\end{document}
